\documentclass[12pt]{article}
\usepackage[a4paper,margin=1in]{geometry}
\usepackage{amsmath, amssymb, booktabs, setspace}
\usepackage{enumitem}
\usepackage{graphicx}
\usepackage{siunitx}

\title{Solutions: Practice Problem Set 5\\ \vspace{10pt} \large ECON 101 -- Macroeconomics}
\author{}
\vspace{-10pt}
\date{ \vspace{-30pt} Winter 2026}

\begin{document}
\maketitle

\begin{enumerate}[label=\textbf{\arabic*.}]

\item \textbf{Three functions of money}

Money performs three main functions:

\textbf{(i) Medium of exchange}\\
Money is used to buy and sell goods and services. It avoids the need for barter.

\textbf{Example:} A student uses cash or a debit balance to buy lunch at the cafeteria.

\textbf{(ii) Unit of account}\\
Money provides a common way of measuring and quoting values.

\textbf{Example:} A textbook is priced at \$120 and a laptop is priced at \$900. These prices are both measured in dollars.

\textbf{(iii) Store of value}\\
Money allows people to transfer purchasing power from the present to the future.

\textbf{Example:} A worker keeps part of her paycheck in a chequing account and uses it next month.

\textbf{Summary:}
\begin{itemize}
    \item Medium of exchange = used for transactions
    \item Unit of account = used to measure value
    \item Store of value = used to hold purchasing power
\end{itemize}

\item \textbf{M1 vs. broader money measures; currency vs. deposits}

\textbf{M1} is a narrow measure of the money supply. It includes the most liquid forms of money, such as:
\begin{itemize}
    \item currency held by the public,
    \item chequable deposits.
\end{itemize}

\textbf{Broader measures of money supply} include M1 plus less liquid assets, such as certain savings accounts and other deposits that can be converted into spending power fairly easily.

\textbf{Currency vs. deposits}
\begin{itemize}
    \item \textbf{Currency} means physical money, such as bank notes and coins.
    \item \textbf{Deposits} are balances held in bank accounts.
\end{itemize}

\textbf{Why are some assets included in the money supply while others are not?}

The key criterion is \textbf{liquidity}. Assets included in the money supply can be used directly or very easily for transactions.

\begin{itemize}
    \item Currency and chequable deposits are highly liquid.
    \item Long-term bonds and some financial assets are not included because they are not used directly to buy goods and services.
\end{itemize}

\textbf{Conclusion:} The more easily an asset can be used for transactions, the more likely it is to be included in the money supply.

\item \textbf{How commercial banks create money}

Commercial banks create money through the process of lending.

\textbf{Step 1: Initial deposit}\\
Suppose a customer deposits money in a bank.

\textbf{Step 2: Reserve requirement / desired reserve ratio}\\
The bank keeps a fraction of the deposit as reserves and lends out the rest.

\textbf{Step 3: New loan becomes new deposit}\\
When the loan is spent, it often gets redeposited into the banking system.

\textbf{Step 4: Repeated lending}\\
The next bank keeps part as reserves and lends the rest again.

This process continues through the banking system, causing total deposits to increase.

\textbf{Role of the reserve ratio}

The reserve ratio determines how much of each deposit banks keep instead of lending.

\begin{itemize}
    \item A \textbf{higher} reserve ratio means banks lend less, so money creation is smaller.
    \item A \textbf{lower} reserve ratio means banks lend more, so money creation is larger.
\end{itemize}

Thus:
\[
\boxed{\text{Banks create money by making loans, and the reserve ratio limits the size of deposit creation.}}
\]

\end{enumerate}

\begin{enumerate}[label=\textbf{\arabic*.}, start=4]

\item \textbf{Deposit Creation}

Given:
\begin{itemize}
    \item Initial deposit = \$5{,}000
    \item Desired reserve ratio \(r_d = 0.10\)
\end{itemize}

\begin{enumerate}
    \item[(a)] \textbf{Calculate required reserves}

Required (or desired) reserves are:
\[
\text{Reserves} = r_d \times \text{deposit}
\]

Substitute:
\[
\text{Reserves} = 0.10 \times 5{,}000 = 500
\]

\[
\boxed{\text{Required reserves} = \$500}
\]

\item[(b)] \textbf{Calculate initial loans}

The bank lends out the remainder of the deposit:
\[
\text{Loans} = \text{deposit} - \text{reserves}
\]

\[
\text{Loans} = 5{,}000 - 500 = 4{,}500
\]

\[
\boxed{\text{Initial loans} = \$4{,}500}
\]

\item[(c)] \textbf{Calculate the simple deposit multiplier}

The simple deposit multiplier is:
\[
m = \frac{1}{r_d}
\]

Substitute:
\[
m = \frac{1}{0.10} = 10
\]

\[
\boxed{m = 10}
\]

\item[(d)] \textbf{Calculate the maximum increase in deposits}

Maximum increase in deposits:
\[
\Delta D = m \times \text{initial deposit}
\]

\[
\Delta D = 10 \times 5{,}000 = 50{,}000
\]

\[
\boxed{\text{Maximum increase in deposits} = \$50{,}000}
\]

\textbf{Interpretation:} If all excess reserves are lent out and all loan proceeds are redeposited, the initial \$5,000 deposit can generate up to \$50,000 of total deposits in the banking system.
\end{enumerate}

\item \textbf{Leakages in the banking system}

Two important leakages reduce the money multiplier:

\textbf{(i) Excess reserves}\\
If banks choose to hold reserves above the desired minimum, they lend out less money. Since fewer loans are made, the deposit expansion process becomes smaller.

\textbf{(ii) Currency held by the public}\\
If households keep some cash instead of redepositing loan proceeds into banks, less money returns to the banking system. This reduces further rounds of lending.

\textbf{Conclusion:}
Both excess reserves and currency leakages make the actual multiplier smaller than the simple multiplier.

\item \textbf{Bank runs}

A \textbf{bank run} occurs when many depositors try to withdraw their money at the same time because they fear the bank may fail.

\textbf{Why bank runs happen:}
\begin{itemize}
    \item Banks keep only a fraction of deposits as reserves.
    \item Most deposited funds are lent out.
    \item If too many people demand cash at once, the bank may not have enough liquid reserves.
\end{itemize}

\textbf{How deposit insurance helps}

Deposit insurance protects depositors up to a certain limit if a bank fails. This reduces panic because people know their money is protected.

\textbf{How the Bank of Canada helps}

The Bank of Canada can support financial stability by:
\begin{itemize}
    \item acting as a lender of last resort,
    \item providing liquidity to the financial system,
    \item helping maintain confidence in the banking system.
\end{itemize}

\textbf{Conclusion:} Deposit insurance and central bank support reduce the likelihood and severity of bank runs by preventing panic.

\end{enumerate}



\begin{enumerate}[label=\textbf{\arabic*.}, start=7]

\item \textbf{Bank of Canada inflation target}

The Bank of Canada targets inflation at \textbf{2\%}, which is the midpoint of a control range of \textbf{1\% to 3\%}.

\textbf{Why inflation targeting is important:}
\begin{itemize}
    \item It helps anchor inflation expectations.
    \item It provides a clear goal for monetary policy.
    \item It supports price stability.
    \item Stable inflation helps households and firms make better long-term plans.
\end{itemize}

\textbf{Conclusion:}
Inflation targeting improves credibility, reduces uncertainty, and promotes macroeconomic stability.


\item \textbf{Monetary transmission mechanism}

The monetary transmission mechanism explains how a change in the policy interest rate affects output and inflation.

\textbf{Step 1: Policy rate changes}\\
The Bank of Canada raises or lowers the target overnight rate.

\textbf{Step 2: Other interest rates change}\\
Short-term market rates and many lending rates move in the same direction.

\textbf{Step 3: Consumption and investment respond}\\
\begin{itemize}
    \item Lower interest rates encourage households to borrow and spend.
    \item Lower interest rates encourage firms to invest more.
\end{itemize}

\textbf{Step 4: Exchange rate may change}\\
A lower Canadian interest rate may reduce demand for Canadian financial assets, causing the Canadian dollar to depreciate.

\textbf{Step 5: Net exports respond}\\
A depreciation makes Canadian exports cheaper and imports more expensive, increasing net exports.

\textbf{Step 6: Aggregate demand changes}\\
Changes in consumption, investment, and net exports shift aggregate demand.

\textbf{Step 7: Output and inflation respond}\\
Higher AD raises output in the short run and may increase inflationary pressure. Lower AD reduces inflationary pressure.

\textbf{Conclusion:} Monetary policy affects the economy indirectly through interest rates, exchange rates, and spending decisions.

\item \textbf{Expansionary vs. contractionary monetary policy}

\textbf{Expansionary monetary policy}
\begin{itemize}
    \item The Bank of Canada lowers the policy interest rate.
    \item This encourages borrowing, spending, and investment.
    \item It is used when the economy is weak, output is below potential, or unemployment is high.
\end{itemize}

\textbf{Contractionary monetary policy}
\begin{itemize}
    \item The Bank of Canada raises the policy interest rate.
    \item This discourages borrowing and spending.
    \item It is used when inflation is too high or the economy is overheating.
\end{itemize}

\textbf{Summary:}
\begin{itemize}
    \item Expansionary policy fights recession.
    \item Contractionary policy fights inflation.
\end{itemize}

\end{enumerate}

\begin{enumerate}[label=\textbf{\arabic*.}, start=10]

\item \textbf{Multiplier and monetary policy}

Given:
\begin{itemize}
    \item \(MPC = 0.8\)
    \item Investment increases by \$40 billion
\end{itemize}

\begin{enumerate}
    \item[(a)] \textbf{Calculate the multiplier}

\[
k = \frac{1}{1-MPC}
\]

\[
k = \frac{1}{1-0.8} = \frac{1}{0.2} = 5
\]

\[
\boxed{k = 5}
\]

\item[(b)] \textbf{Calculate the total change in GDP}

\[
\Delta Y = k \times \Delta I
\]

\[
\Delta Y = 5 \times 40 = 200
\]

\[
\boxed{\Delta Y = \$200 \text{ billion}}
\]
\end{enumerate}

\item \textbf{Effects of an increase in interest rates}

\textbf{Consumption:}\\
Higher interest rates make borrowing more expensive, so households reduce spending on consumer durables and housing.

\textbf{Investment:}\\
Higher interest rates increase the cost of borrowing for firms, so investment spending falls.

\textbf{Exchange rate:}\\
Higher Canadian interest rates attract foreign investors, increasing demand for the Canadian dollar. The dollar appreciates.

\textbf{Net exports:}\\
An appreciation makes exports more expensive and imports cheaper, so net exports fall.

\textbf{Conclusion:}
\[
\boxed{\text{Higher interest rates reduce C, I, and NX, which lowers aggregate demand.}}
\]

\item \textbf{Why monetary policy may be less effective during a recession}

Monetary policy may be less effective in a recession for several reasons.

\textbf{1. Liquidity trap / very low interest rates}\\
If interest rates are already very low, the central bank may have little room to reduce them further.

\textbf{2. Weak expectations}\\
Even if rates fall, households and firms may still be pessimistic about the future. They may not borrow or spend more.

\textbf{3. Banking weakness / credit frictions}\\
Banks may be unwilling to lend during a recession, even if the central bank lowers rates.

\textbf{4. Slow transmission}\\
Monetary policy takes time to affect spending, output, and inflation.

\textbf{Conclusion:}

During recessions, low confidence and financial frictions can weaken the effect of monetary policy.


\end{enumerate}



\begin{enumerate}[label=\textbf{\arabic*.}, start=13]

\item \textbf{Discretionary fiscal policy vs. automatic stabilizers}

\textbf{Discretionary fiscal policy} refers to deliberate changes in government spending or taxes made by policymakers.

\textbf{Examples:}
\begin{itemize}
    \item a new infrastructure program,
    \item a temporary tax cut.
\end{itemize}

\textbf{Automatic stabilizers} are features of the tax and transfer system that automatically reduce the size of fluctuations in GDP without new legislation.

\textbf{Examples:}
\begin{itemize}
    \item unemployment insurance,
    \item progressive income taxes.
\end{itemize}

\textbf{Difference:}
\begin{itemize}
    \item Discretionary policy requires a policy decision.
    \item Automatic stabilizers operate automatically.
\end{itemize}

\item \textbf{How government spending and taxes affect aggregate demand}

\textbf{Government spending:}\\
An increase in government spending directly raises aggregate demand because government purchases are a component of AD.

\textbf{Taxes:}\\
A decrease in taxes raises disposable income, which increases household consumption and shifts AD right. An increase in taxes does the opposite.

\textbf{Conclusion:}
\[
\boxed{\text{Higher G or lower taxes increase AD; lower G or higher taxes decrease AD.}}
\]

\item \textbf{Budget deficit and surplus}

A \textbf{budget deficit} occurs when government spending exceeds tax revenue:
\[
G > T
\]

A \textbf{budget surplus} occurs when tax revenue exceeds government spending:
\[
T > G
\]

\textbf{Relationship to the business cycle:}
\begin{itemize}
    \item During recessions, tax revenue tends to fall and transfer payments rise, so deficits often increase.
    \item During expansions, tax revenue tends to rise and transfer payments fall, so deficits tend to shrink or surpluses may appear.
\end{itemize}

Thus, part of the budget balance changes automatically with the business cycle.

\end{enumerate}

\begin{enumerate}[label=\textbf{\arabic*.}, start=16]

\item \textbf{Government Spending Multiplier}

Given:
\begin{itemize}
    \item \(MPC = 0.75\)
    \item Government spending increases by \$80 billion
\end{itemize}

\begin{enumerate}
    \item[(a)] \textbf{Calculate the multiplier}

\[
k = \frac{1}{1-MPC}
\]

\[
k = \frac{1}{1-0.75} = \frac{1}{0.25} = 4
\]

\[
\boxed{k = 4}
\]

\item[(b)] \textbf{Calculate the change in GDP}

\[
\Delta Y = k \times \Delta G
\]

\[
\Delta Y = 4 \times 80 = 320
\]

\[
\boxed{\Delta Y = \$320 \text{ billion}}
\]
\end{enumerate}

\item \textbf{Tax Multiplier}

Given:
\begin{itemize}
    \item \(MPC = 0.75\)
    \item Taxes decrease by \$60 billion
\end{itemize}

\begin{enumerate}
    \item[(a)] \textbf{Calculate the tax multiplier}

The tax multiplier is:
\[
k_T = -\frac{MPC}{1-MPC}
\]

Substitute:
\[
k_T = -\frac{0.75}{1-0.75} = -\frac{0.75}{0.25} = -3
\]

\[
\boxed{k_T = -3}
\]

\item[(b)] \textbf{Calculate the change in GDP}

Since taxes \emph{decrease} by 60, we have:
\[
\Delta T = -60
\]

Then:
\[
\Delta Y = k_T \times \Delta T
\]

\[
\Delta Y = (-3)\times(-60)=180
\]

\[
\boxed{\Delta Y = \$180 \text{ billion}}
\]
\end{enumerate}

\end{enumerate}

\begin{enumerate}[label=\textbf{\arabic*.}, start=18]

\item \textbf{Fiscal vs. monetary policy in a recession}

Given:
\begin{itemize}
    \item GDP is below potential
    \item unemployment is rising
\end{itemize}

\begin{enumerate}
    \item[(a)] \textbf{Use the AE model to explain the recession}

In the AE model, a recession occurs when planned aggregate expenditure is too low. Firms find that sales are weak, inventories rise unexpectedly, and they cut production. As output falls, income falls, which further weakens spending through the multiplier process.

\item[(b)] \textbf{Use the AD--AS model to illustrate the situation}

In the AD--AS model:
\begin{itemize}
    \item Aggregate demand is too low.
    \item The AD curve lies to the left of the level needed for full-employment output.
    \item Equilibrium output is below potential GDP.
    \item Unemployment rises.
\end{itemize}

\item[(c)] \textbf{One fiscal policy and one monetary policy}

\textbf{Fiscal policy:} Increase government spending or decrease taxes.\\
This shifts AE upward and AD to the right, raising GDP.

\textbf{Monetary policy:} Lower the policy interest rate.\\
This encourages consumption, investment, and net exports, shifting AD to the right.

\textbf{Conclusion:} Both policies aim to increase aggregate demand and close the recessionary gap.
\end{enumerate}

\end{enumerate}

\begin{enumerate}[label=\textbf{\arabic*.}, start=19]

\item \textbf{Crowding out}

Suppose the government increases spending by \$100 billion and finances it through borrowing.

\textbf{Step 1: Government borrowing rises}\\
A larger deficit means the government demands more loanable funds.

\textbf{Step 2: Interest rates rise}\\
With higher demand for loanable funds, the equilibrium interest rate tends to increase.

\textbf{Step 3: Private investment falls}\\
Higher interest rates make borrowing more expensive for firms, so private investment spending decreases.

\textbf{Step 4: Actual GDP rises by less than predicted}\\
The simple multiplier assumes no offsetting effects. But if private investment falls, part of the increase in government spending is offset.

This is called \textbf{crowding out}.

\textbf{Conclusion:}
Borrowing-financed government spending may raise GDP by less than the simple multiplier predicts because higher interest rates reduce private investment.


\end{enumerate}

\hrule
\vspace{1em}

\begin{center}
\textit{End of Solution Key}
\end{center}

\end{document}